% Source: physical PDF pp. 444--445 (printed pp. 421--422). Coverage:
% Chapter 23 opening, from the chapter heading through immediately before
% Section 23.1; citations [1]--[7], no equations, figures, tables, or
% footnotes. Corrected erratum: the duplicated “as” before “magnetic monopoles” is removed.

Most of this book has been devoted to applications of quantum field theory
that can at least be described in perturbation theory, whether or not the
perturbation series actually works well numerically. In using perturbation
theory, we expand the action around the usual spacetime-independent vacuum
values of the fields, keeping the leading quadratic term in the exponential
$\exp(\ii I)$, and treating all terms of higher order in the fields as small
corrections. Starting in the mid-1970s, there has been a growing interest in
effects that arise because there are extended spacetime-dependent field
configurations, such as those known as instantons%
~\hyperref[ch23-ref-1]{[1]}, that are also stationary `points' of the
action. In principle, we must include these configurations in path integrals
and sum over fluctuations around them. (In
\hyperref[sec:20.7]{Section 20.7} we have already seen an example of an
instanton configuration, applied in a different context.) Although such
non-perturbative contributions are often highly suppressed, they are large
in quantum chromodynamics, and produce interesting exotic effects in the
standard electroweak theory.

There are also extended field configurations that occur, not only as
correction terms in path integrals for processes involving ordinary
particles, but also as possible components of actual physical states. These
configurations include some that are particle-like, such as magnetic
monopoles~\hyperref[ch23-ref-2]{[2]} and
skyrmions~\hyperref[ch23-ref-3]{[3]}, which are concentrated around a point
in space or, equivalently, around a world line in spacetime. There are also
string-like configurations~\hyperref[ch23-ref-4]{[4]}, similar to the
vortex lines in superconductors discussed in
\hyperref[sec:21.6]{Section 21.6}, which are concentrated around a line in
space or, equivalently, around a world sheet in spacetime. Then there are
configurations that are sheet-like, like the domain
walls~\hyperref[ch23-ref-5]{[5]} between spatial regions in which discrete
symmetries are broken in different ways. In contrast, the instantons
mentioned above are \emph{event-like}, concentrated about a point in
spacetime, and therefore never appear as components of actual physical
states.

Some extended field configurations are stabilized because of boundary
conditions that are imposed by the nature of the problem in which they
appear. An example is the `bounce' solution, which appears in the analysis
of vacuum decay~\hyperref[ch23-ref-6]{[6]}, and will be discussed in
\hyperref[sec:23.8]{Section 23.8}. Other configurations are stable because
they carry a quantum number whose conservation forbids any possible decay
mode~\hyperref[ch23-ref-7]{[7]}.

In this chapter we shall mostly be concerned with extended field
configurations that are stabilized by their topology. In analyzing all such
configurations we use the same topological tools, chiefly homotopy theory,
so we shall begin by considering all topologically stabilized
configurations together, in a space or spacetime of arbitrary dimensionality
$d$.
